\documentclass{article}
\title{Prime Spacing from the C(r) Kernel}
\author{Brian Tice Sr. \\ Axion Resonance Labs}
\date{November 2026}

\begin{document}

\maketitle

\section{Claim}

The average gap structure of the prime numbers is not an arithmetic accident. It is produced by the polynomial growth term in the connectivity kernel

\[
C(\rho) = (1 + 2\rho) e^{-\rho/3} e^{i \pi \rho / 4}
\]

evaluated at the anti-resonance nodes of the lattice. No external number-theoretic input is required.

\section{Anti-resonance nodes}

An anti-resonance node is a lattice site at which the transfer does not couple. Index these nodes by $k = 1, 2, 3, \dots$. At node $k$ the local coupling capacity contributed by the polynomial factor is

\[
c_k = \frac{1}{1 + 2k}.
\]

\section{Cumulative capacity}

The cumulative coupling capacity up to node $n$ is the partial sum

\[
S_n = \sum_{k=1}^{n} \frac{1}{1 + 2k} = \frac{1}{2} \sum_{k=1}^{n} \frac{1}{k + 1/2}.
\]

This is a shifted harmonic series. Its integral approximation is

\[
S_n \approx \frac{1}{2} \log n + C,
\]

where $C$ is a constant (the Euler--Mascheroni shift for half-integers).

\section{Threshold crossings}

Anti-resonance events fire when cumulative capacity crosses successive integer thresholds $m = 1, 2, 3, \dots$. The node index $n_m$ at which threshold $m$ is crossed therefore satisfies

\[
S_{n_m} \approx m \implies \frac{1}{2} \log n_m \approx m - C \implies n_m \approx e^{2(m - C)}.
\]

\section{Gap structure}

The gap between consecutive crossing nodes is

\[
g_m = n_{m+1} - n_m \approx e^{2(m+1-C)} - e^{2(m-C)} = e^{2(m-C)} \bigl( e^{2} - 1 \bigr).
\]

Since $n_m \sim e^{2m}$, the gap scales as

\[
g_m \sim \kappa \, n_m, \qquad \kappa = e^{2} - 1 \approx 6.389.
\]

In terms of the node index itself the average gap is proportional to $n$, which in the continuum limit is the statement that average gaps grow as $\log N$ when $N$ counts all integers up to the node. This is the prime-gap law: the $n$th prime satisfies $p_n \sim n \log n$, so successive gaps average $\log n$.

\section{What is not imported}

\begin{itemize}
    \item No prime number theorem.
    \item No Euler product.
    \item No zeta function.
    \item No fitted constants. The factor $1 + 2\rho$ is the kernel's own polynomial term; the exponential decay $e^{-\rho/3}$ sets the scale at which attenuation dominates but does not alter the gap law.
\end{itemize}

\section{Falsifier}

If the kernel polynomial were replaced by any term whose partial sums grow slower than logarithmically (e.g., bounded, or growing as $\log \log n$), the derived gap law would fall below $\log n$ and the match to prime spacing would fail. The harmonic growth of $1 + 2\rho$ is necessary.

\section{Status}

Derivation complete inside the C(r) framework. Independent numerical check: compare the sequence of threshold-crossing nodes $n_m$ against the sequence of primes and measure the correlation of gap distributions.

\end{document}